\makeatletter

% ----------------------------------------------------------------------------
% load-first.tex

\usepackage{amsmath}
\usepackage{amsthm}
\usepackage{mathtools}
\mathtoolsset{
    showonlyrefs=true,
    showmanualtags=true,
}

\usepackage{iftex}

% ----------------------------------------------------------------------------
% colors.tex

\usepackage[dvipsnames, table]{xcolor}

\definecolor{CardinalRed}{HTML}{8C1515}
\definecolor{StanfordWhite}{HTML}{FFFFFF}
\definecolor{StanfordBlack}{HTML}{2E2D29}
\definecolor{CoolGrey}{HTML}{53565A}
\definecolor{PaloAlto}{HTML}{175E54}
\definecolor{PaloVerde}{HTML}{279989}
\definecolor{Olive}{HTML}{8F993E}
\definecolor{Bay}{HTML}{6FA287}
\definecolor{Sky}{HTML}{4298B5}
\definecolor{Lagunita}{HTML}{007C92}
\definecolor{Poppy}{HTML}{E98300}
\definecolor{Spirited}{HTML}{E04F39}
\definecolor{Illuminating}{HTML}{FEDD5C}
\definecolor{Plum}{HTML}{620059}
\definecolor{Brick}{HTML}{651C32}
\definecolor{Archway}{HTML}{5D4B3C}
\definecolor{Stone}{HTML}{7F7776}
\definecolor{Fog}{HTML}{DAD7CB}

% ----------------------------------------------------------------------------
% fonts.tex

\newcommand{\firaSansScale}{0.93}
\newcommand{\firaMonoScale}{0.93}

\ifluatex
    \usepackage{fontspec}
    \usepackage[
        math-style=ISO, bold-style=ISO,
        partial=upright,
        mathrm=sym, mathit=sym, mathsf=sym, mathbf=sym, mathtt=sym,
        warnings-off={mathtools-colon,mathtools-overbracket},
    ]{unicode-math}

    \setmainfont{TeX Gyre Pagella}
    \setsansfont{Fira Sans}[Scale=\firaSansScale]
    \IfFontExistsTF{JetBrains Mono}{
        \setmonofont{JetBrains Mono}[Scale=MatchLowercase]
    }{
        \setmonofont{Fira Mono}[Scale=\firaMonoScale]
    }

    \setmathfont{TeX Gyre Pagella Math}
    \setmathfont{NewCMMath-Book.otf}[
        range={
                \mathopen,\mathclose,\mathfence,\mathop,\mathbin,\mathrel,
                \mathcal,\varnothing
            },
    ]
    \setmathfont{TeX Gyre Pagella Math}[range={\mathexclam}]
    \setmathfont{NewCMMath-Book.otf}[range={}]
\else
    \usepackage{newpx}
    \usepackage[scaled=\firaSansScale]{FiraSans}
    \usepackage[scaled=\firaMonoScale]{FiraMono}
\fi

% ----------------------------------------------------------------------------
% links.tex

\usepackage{zref-clever}

\DeclareDocumentCommand\cref{ s m }{%
    \IfBooleanTF{#1}{%
        \zcref*[nameinlink=false, cap]{#2}%
    }{%
        \zcref[nameinlink=false, cap]{#2}%
    }%
}

% ----------------------------------------------------------------------------
% theorems.tex

\newtheoremstyle{it}{}{}{\itshape}{}{\bfseries}{}{0.5em}{%
    \thmname{#1}\thmnumber{ #2}\thmnote{ (#3)}.}

\newtheoremstyle{rm}{}{}{\rmfamily}{}{\bfseries}{}{0.5em}{%
    \thmname{#1}\thmnumber{ #2}\thmnote{ (#3)}.}

\providecommand{\theoremTheoremStyle}{it}
\providecommand{\exampleTheoremStyle}{rm}

\newcommand{\newTheoremWithOwnCounter}[2]{
    \newtheorem{#1}{#2}
    \newtheorem*{#1*}{#2}
    \zcRefTypeSetup{#1}{name-sg=#1,Name-sg=#2,name-pl=#1s,Name-pl=#2s}
}
\newcommand{\newTheoremWithSharedCounter}[3]{%
    \newtheorem{#1}[#3]{#2}
    \newtheorem*{#1*}{#2}
    \zcRefTypeSetup{#1}{name-sg=#1,Name-sg=#2,name-pl=#1s,Name-pl=#2s}
    \AddToHook{env/#1/begin}{\zcsetup{countertype={#3=#1}}}
}

\theoremstyle{\theoremTheoremStyle}
\newTheoremWithOwnCounter{theorem}{Theorem}
\newTheoremWithSharedCounter{proposition}{Proposition}{theorem}
\newTheoremWithSharedCounter{lemma}{Lemma}{theorem}
\newTheoremWithSharedCounter{corollary}{Corollary}{theorem}
\newTheoremWithSharedCounter{conjecture}{Conjecture}{theorem}
\newTheoremWithOwnCounter{assumption}{Assumption}
\newTheoremWithOwnCounter{claim}{Claim}
\newTheoremWithOwnCounter{definition}{Definition}
\newTheoremWithOwnCounter{remark}{Remark}

\theoremstyle{\exampleTheoremStyle}
\newTheoremWithOwnCounter{example}{Example}
\newTheoremWithOwnCounter{exercise}{Exercise}

\newenvironment{exampleQed}{%
    \renewcommand{\qedsymbol}{\( \triangleleft \)}%
    \pushQED{\qed}\example
}{\popQED\endexample}

\renewenvironment{proof}[1][\proofname]{\par
    \pushQED{\qed}%
    \normalfont \topsep6\p@\@plus6\p@\relax
    \trivlist
    \item[\hskip\labelsep\bfseries #1\@addpunct{.}]\ignorespaces
}{%
    \popQED\endtrivlist\@endpefalse
}

\newenvironment{subproof}[1][\proofname]{%
    \renewcommand{\qedsymbol}{\( \triangleleft \)}%
    \begin{proof}[\sc #1]%
}{%
    \end{proof}%
}

% ----------------------------------------------------------------------------
% math.tex

\usepackage{mleftright}
\mleftright{}

\usepackage{bracealign}

\newcommand{\opn}[1]{\operatorname{#1}}

\newcommand{\R}{\mathbb{R}}
\newcommand{\Z}{\mathbb{Z}}

\newcommand{\N}{\mathcal{N}}

\newcommand{\e}{\mathrm{e}}
\newcommand{\ii}{\mathrm{i}}
\newcommand{\ppi}{\mathrm{\pi}}

\newcommand{\bmid}{\mathbin{\mid}}

\DeclareMathOperator{\E}{\mathbb{E}}
\DeclareMathOperator{\Var}{Var}
\DeclareMathOperator{\Cov}{Cov}

\DeclareMathOperator*{\argmax}{arg\,max}
\DeclareMathOperator*{\argmin}{arg\,min}
\DeclareMathOperator*{\essinf}{ess\,inf}
\DeclareMathOperator*{\esssup}{ess\,sup}
\DeclareMathOperator*{\maximize}{maximize}
\DeclareMathOperator*{\minimize}{minimize}

\DeclareDocumentCommand{\ol}{ m }{
    {\mspace{1mu}\overline{\mspace{-1mu} #1 \mspace{-1mu}}\mspace{1mu}}
}
\DeclareDocumentCommand{\ul}{ m }{
    {\mspace{1mu}\underline{\mspace{-1mu} #1 \mspace{-1mu}}\mspace{1mu}}
}

\def\cases{\left\{\def\arraystretch{1.2}\;\hskip-\arraycolsep \array{l@{\quad}l}}
\def\endcases{\endarray\hskip-\arraycolsep\right.}

\renewcases{dcases}{\quad}{$\m@th\displaystyle##$\hfil}{$\m@th\displaystyle##$\hfil}{\lbrace\;}{.}

\DeclareDocumentCommand{\abs}{ s m }{
    \IfBooleanTF{#1}{
        \lvert #2 \rvert
    }{
        \left\lvert #2 \right\rvert
    }
}
\DeclareDocumentCommand{\norm}{ s m }{
    \IfBooleanTF{#1}{
        \lVert #2 \rVert
    }{
        \left\lVert #2 \right\rVert
    }
}

\DeclareDocumentCommand{\qq}{ s m }{
    \IfBooleanTF{#1}{}{\quad} \text{#2} \quad
}

\usepackage{fixdif}
\letdif{\grad}{nabla}

\DeclareDocumentCommand\flatfrac{ m m }{ #1 / #2 }

\DeclareDocumentCommand\dv{ s o m g }{
    \IfBooleanTF{#1}{ \let\fractype\flatfrac }{ \let\fractype\frac }
    \IfNoValueTF{#4}{
        \fractype{
            \d \IfNoValueTF{#2}{}{^{#2}}
        }{
            \d #3\IfNoValueTF{#2}{}{^{#2}}
        }
    }{
        \fractype{
            \d \IfNoValueTF{#2}{}{^{#2}} #3
        }{
            \d #4\IfNoValueTF{#2}{}{^{#2}}
        }
    }
}

\DeclareDocumentCommand\pdv{ s o m g }{
    \IfBooleanTF{#1}{ \let\fractype\flatfrac }{ \let\fractype\frac }
    \IfNoValueTF{#4}{
        \fractype{
            \partial \IfNoValueTF{#2}{}{^{#2}}
        }{
            \partial #3\IfNoValueTF{#2}{}{^{#2}}
        }
    }{
        \fractype{
            \partial \IfNoValueTF{#2}{}{^{#2}} #3
        }{
            \partial #4\IfNoValueTF{#2}{}{^{#2}}
        }
    }
}

% ----------------------------------------------------------------------------
% lists.tex

\@ifclassloaded{beamer}{ }{
    \usepackage{enumitem}

    \setlist*[itemize]{label=\textbullet}

    \newlist{cenumerate}{enumerate}{4}
    \newlist{citemize}{itemize}{4}
    \setlist*[cenumerate,citemize,1]{noitemsep}
    \setlist*[cenumerate,citemize,2,3,4]{nosep}
    \setlist*[citemize]{label=\textbullet}
    \setlist*[cenumerate,1]{label=\arabic*.}
    \setlist*[cenumerate,2]{label=\alph*)}
    \setlist*[cenumerate,3]{label=\roman*)}
    \setlist*[cenumerate,4]{label=\Alph*)}
}

% ----------------------------------------------------------------------------
% misc.tex

\@ifclassloaded{article}{
    \usepackage[a4paper, vscale=0.75]{geometry}
}

\usepackage{setspace}
\setstretch{1.2}

% No first-paragraph indentation for quotations and abstracts
\appto\quotation{\noindent\ignorespaces}

\usepackage{caption}

\usepackage{tabularray}
\UseTblrLibrary{booktabs}

\usepackage[maxlevel=4]{csquotes}

\usepackage[authordate, backend=biber]{biblatex-chicago}

\usepackage{microtype}

\usepackage{ellipsis}

\usepackage{url}
\g@addto@macro{\UrlBreaks}{\UrlOrds}

\usepackage{silence}
\WarningFilter{microtype}{Unable to apply patch `footnote'}

\usepackage[]{hyperref}
\hypersetup{
    allcolors = .,
    breaklinks = true,
    colorlinks = true,
    hypertexnames = false,
    bookmarksnumbered = true,
}
\urlstyle{rm}

\renewcommand*{\@seccntformat}[1]{%
    \csname the#1\endcsname.\quad
}
\renewcommand\Hy@numberline[1]{#1. }

\let\oldFootnote\footnote
\newcommand\nextToken\relax
\renewcommand\footnote[1]{\oldFootnote{#1}\futurelet\nextToken\isFootnote}
\newcommand\isFootnote{\ifx\footnote\nextToken\textsuperscript{,}\fi}

\NewDocumentCommand{\TwoPartSection}{o m m}{%
    \IfNoValueTF{#1}
    {\section[#2]{#2\\\textmd{\textit{\large #3}}}}%
    {\section[#1]{#2\\\textmd{\textit{\large #3}}}}%
}
\NewDocumentCommand{\TwoPartSubsection}{o m m}{%
    \IfNoValueTF{#1}
    {\subsection[#2]{#2\\\textmd{\textit{\large #3}}}}%
    {\subsection[#1]{#2\\\textmd{\textit{\large #3}}}}%
}

% ----------------------------------------------------------------------------
% beamer.tex

\@ifclassloaded{beamer}{
    \usetheme{moloch}

    \molochcolors{theme=paper}
    \molochcolors{light/title separator=CardinalRed}
    \hypersetup{allcolors = .}

    \setlength{\parskip}{1em plus 0.2em minus 0.2em}
    \newcommand{\reduceSpacingIfZeroDepth}{\ifnum\@itemdepth=0 \vspace{-1em}\fi}
    \AddToHook{env/itemize/before}{\reduceSpacingIfZeroDepth}
    \AddToHook{env/enumerate/before}{\reduceSpacingIfZeroDepth}

    \molochset{titleformat frame=smallcaps}
    \setbeamertemplate{page number in head/foot}[totalframenumber]

    \ExecuteBibliographyOptions{maxcitenames=2}
    \setbeamertemplate{bibliography item}[text]
}{}

% ----------------------------------------------------------------------------
% project.tex

\makeatother
